\documentclass[aspectratio=169,10pt]{beamer}
\usetheme[numbering=fraction,progressbar=frametitle]{metropolis}
\usepackage{fontspec}
\setsansfont{Latin Modern Sans}
\usepackage{amsmath,amssymb,mathtools,graphicx,booktabs}
\usepackage{polyglossia}
\setdefaultlanguage{french}
\setbeamertemplate{navigation symbols}{}
\setbeamerfont{frametitle}{size=\large}
\setbeamerfont{normal text}{size=\normalsize}
\setlength{\parskip}{.45em}
\hypersetup{pdfauthor={},pdfsubject={Lecture slides},pdfkeywords={},pdftitle={Orthogonalité et distances}}
\providecommand{\N}{\mathbb{N}}
\providecommand{\Z}{\mathbb{Z}}
\providecommand{\D}{\mathbb{D}}
\providecommand{\Q}{\mathbb{Q}}
\providecommand{\R}{\mathbb{R}}
\providecommand{\C}{\mathbb{C}}
\providecommand{\vect}[1]{\overrightarrow{#1}}
\providecommand{\vi}{\vec{\imath}}
\providecommand{\vj}{\vec{\jmath}}
\providecommand{\vk}{\vec{k}}
\providecommand{\norm}[1]{\left\lVert #1\right\rVert}
\providecommand{\vc}[2]{\begin{pmatrix}#1\\#2\end{pmatrix}}
\providecommand{\deter}[2]{\det\left(#1,#2\right)}
\providecommand{\abs}[1]{\left\lvert #1\right\rvert}
\providecommand{\dd}{\,\mathrm{d}}
\providecommand{\e}{\mathrm{e}}
\providecommand{\Card}{\operatorname{Card}}
\providecommand{\E}{\operatorname{E}}
\providecommand{\V}{\operatorname{V}}
\providecommand{\degre}{^{\circ}}
\providecommand{\un}[1]{\,\mathrm{#1}}
\providecommand{\sci}[2]{#1\times10^{#2}}
\providecommand{\pgcd}{\operatorname{PGCD}}
\providecommand{\ppcm}{\operatorname{PPCM}}
\title{Orthogonalité et distances}
\subtitle{Cours et exemples}
\date{}
\author{}
\begin{document}
\begin{frame}[plain]
\vfill
{\usebeamerfont{title}\inserttitle\par}
\medskip
{\usebeamerfont{subtitle}\insertsubtitle\par}
\medskip
{\usebeamercolor[fg]{progress bar}\rule{\textwidth}{.5pt}}
\vfill
\end{frame}
\begin{frame}{Objectifs}
Calculer un produit scalaire ; justifier une orthogonalité ; déterminer un projeté et une distance ; calculer un volume.
\end{frame}
\begin{frame}{Notes de cours}
Les calculs par coordonnées de ce chapitre se font dans un repère orthonormé de l'espace.
\end{frame}
\begin{frame}{Produit scalaire dans l'espace}
Pour des vecteurs non nuls \(\vec u,\vec v\), \(\vec u\cdot\vec v=\|\vec u\|\|\vec v\|\cos\theta\), où \(\theta\in[0,\pi]\) est leur angle. Si l'un est nul, le produit scalaire vaut \(0\).
\[\vec u\cdot\vec v=\tfrac12(\|\vec u+\vec v\|^2-\|\vec u\|^2-\|\vec v\|^2).
\]
Dans un repère orthonormé, pour \(\vec u=(x;y;z)\) et \(\vec v=(x';y';z')\),
\[\vec u\cdot\vec v=xx'+yy'+zz'.\]
Le produit scalaire est symétrique et bilinéaire ; \(\vec u\cdot\vec u=\|\vec u\|^2\). Deux vecteurs sont orthogonaux si et seulement si leur produit scalaire est nul.
\end{frame}
\begin{frame}{Produit scalaire dans l'espace}
\(\vec u=(1;2;2)\) et \(\vec v=(2;-1;0)\) sont orthogonaux : \(2-2+0=0\). Leurs normes valent \(3\) et \(\sqrt5\).
\end{frame}
\begin{frame}{Droite, plan et vecteur normal}
Deux droites sont orthogonales si leurs vecteurs directeurs le sont. Une droite est orthogonale à un plan si elle est orthogonale à deux directions non colinéaires de ce plan. Un vecteur normal à un plan est un vecteur non nul orthogonal à toutes ses directions.
\end{frame}
\begin{frame}{Droite, plan et vecteur normal}
Si \(\vec n\cdot\vec u=\vec n\cdot\vec v=0\), alors, pour tout \(\vec w=s\vec u+t\vec v\), \(\vec n\cdot\vec w=s\vec n\cdot\vec u+t\vec n\cdot\vec v=0\).
\end{frame}
\begin{frame}{Droite, plan et vecteur normal}
Si \(A\) appartient au plan \(\mathcal P\) de vecteur normal \(\vec n\), alors \(M\in\mathcal P\) si et seulement si \(\vect{AM}\cdot\vec n=0\).
\end{frame}
\begin{frame}{Projection orthogonale et distance}
Le projeté orthogonal \(H\) de \(M\) sur un plan est l'unique point du plan tel que \(\vect{MH}\) soit normal au plan ; si \(M\) appartient au plan, \(H=M\). La distance de \(M\) au plan est \(MH\). Sur une droite de directeur \(\vec u\), le projeté vérifie \(H\in d\) et \(\vect{MH}\cdot\vec u=0\).
\end{frame}
\begin{frame}{Projection orthogonale et distance}
Pour tout point \(P\) du plan, \(\vect{MH}\perp\vect{HP}\), donc \(MP^2=MH^2+HP^2\ge MH^2\). L'égalité a lieu seulement pour \(P=H\). Le projeté est le point du plan le plus proche.
\end{frame}
\begin{frame}{Projection orthogonale et distance}
Si \(d=A+\R\vec u\), écrire \(H=A+t\vec u\) et résoudre \((A-M+t\vec u)\cdot\vec u=0\). Ainsi \(t=\dfrac{\vect{AM}\cdot\vec{u}}{\lVert\vec u\rVert^2}\).
\end{frame}
\begin{frame}{Angles, aires et volumes}
Pour un angle \(\widehat{ABC}\), utiliser \(\vect{BA}\) et \(\vect{BC}\) : \(\cos \widehat{ABC}=\dfrac{\vect{BA}\cdot\vect{BC}}{BA\times BC}\). Pour l'angle aigu entre deux droites, prendre la valeur absolue du quotient.
L'aire d'un triangle est \(\dfrac{bh}{2}\) ; le volume d'un tétraèdre est \(\dfrac{\mathcal{A}_{\rm base}h}{3}\), où \(h\) est une distance orthogonale. Les unités sont respectivement au carré et au cube.
\end{frame}
\begin{frame}{Développements et lieux géométriques}
Par bilinéarité,
\[\|\vec u+\vec v\|^2=\|\vec u\|^2+2\vec u\cdot\vec v+\|\vec v\|^2,\qquad
\|\vec u-\vec v\|^2=\|\vec u\|^2-2\vec u\cdot\vec v+\|\vec v\|^2.\]
On peut donc aussi calculer \(\vec u\cdot\vec v=\dfrac12(\|\vec u\|^2+\|\vec v\|^2-\|\vec u-\vec v\|^2)\).
Pour deux points distincts \(A,B\) de milieu \(I\), le lieu des points \(M\) tels que \(MA=MB\) est le plan passant par \(I\) de normal \(\vect{AB}\) : le \emph{plan médiateur}. En effet, en écrivant les vecteurs depuis \(I\), les développements donnent \(MA^2-MB^2=2\vect{IM}\cdot\vect{AB}\). L'égalité des distances équivaut donc à \(\vect{IM}\cdot\vect{AB}=0\).
\textbf{Complément, sans évaluation dédiée en 2026-2027.} Deux plans sont perpendiculaires lorsque leurs vecteurs normaux sont orthogonaux. Il ne faut pas confondre cette propriété avec le parallélisme des plans, qui correspond à des normaux colinéaires.
\end{frame}
\begin{frame}{Activité de recherche}
\textbf{Le point d'une droite le plus proche - 35 min.}

Dans un repère orthonormé de l'espace, on considère \(A(1;0;0)\), \(M(0;2;3)\) et la droite \(d\) passant par \(A\), de vecteur directeur \(\vec u=(1;1;0)\). On cherche le point de \(d\) le plus proche de \(M\).
\end{frame}
\begin{frame}{Activité de recherche (2)}
\begin{enumerate}
\setcounter{enumi}{0}\item Exprimer, en fonction d'un paramètre réel \(t\), les coordonnées d'un point \(H\) de \(d\).

\item Déterminer la valeur de \(t\) pour laquelle \(\vect{MH}\cdot\vec u=0\). En déduire les coordonnées du point \(H\) ainsi obtenu.

\item Calculer la distance \(MH\).

\item Soit \(P\) un point quelconque de \(d\). Démontrer que \(MP^2=MH^2+HP^2\), puis en déduire que \(H\) est l'unique point de \(d\) le plus proche de \(M\).

\end{enumerate}
\end{frame}
\end{document}
